% Generated draft: requires mathematical and source review.
\documentclass[11pt]{article}
\usepackage[margin=25mm]{geometry}
\usepackage{fontspec}
\setmainfont{Latin Modern Roman}
\usepackage{amsmath,amssymb,mathtools,unicode-math}
\setmathfont{Latin Modern Math}
\usepackage{enumitem,needspace,xurl,hyperref}
\setlength{\parindent}{0pt}
\setlength{\parskip}{5pt}
\setlength{\emergencystretch}{3em}
\newcommand{\N}{\mathbb{N}}
\newcommand{\Z}{\mathbb{Z}}
\newcommand{\Q}{\mathbb{Q}}
\newcommand{\R}{\mathbb{R}}
\newcommand{\C}{\mathbb{C}}
\newcommand{\F}{\mathbb{F}}
\newcommand{\A}{\mathbb{A}}
\newcommand{\PP}{\mathbb{P}}
\newcommand{\E}{\mathbb{E}}
\newcommand{\eps}{\varepsilon}
\newcommand{\dd}{\,\mathrm{d}}
\begin{document}
\section*{UnsolvedMath ID: 30004953}
\textbf{Source problem code:} catalogue entry
\subsection{Definitions and mathematical statement}
Let \(p\in\R\) satisfy \(p\leq -1\). Write \(S^{n-1}=\{x\in\R^n:\|x\|_2=1\}\). A proper great subsphere has the form \(\xi_L=S^{n-1}\cap L\), where \(L\) is a nonzero proper linear subspace of \(\R^n\). The supplied record does not specify the allowed ambient dimensions or whether every dimension of proper great subsphere is intended. Whenever \(\mu\) is a measure defined on these sets and \(0<\mu(S^{n-1})<\infty\), define its concentration on \(\xi_L\) by
\[
 r_\mu(L)=\frac{\mu(S^{n-1}\cap L)}{\mu(S^{n-1})}.
\]
The positivity and finiteness condition here makes the displayed ratio meaningful; the full measure hypotheses of the cited theorem are not supplied.

\textbf{Problem as supplied.} For \(p\leq -1\), determine the optimal upper restrictions on \(r_\mu(L)\), over the intended great subspheres, that guarantee the existence of a convex body whose \(L_p\) integral curvature measure is \(\mu\). In particular, the original passage asks whether the measure concentration bound in Theorem~2 is optimal.

For notation only, if \(J_p(K,\cdot)\) denotes the intended \(L_p\) integral curvature measure, the requested sufficient condition is to imply
\[
 \exists\,K\quad J_p(K,\cdot)=\mu
\]
for every measure \(\mu\) satisfying the intended hypotheses. Neither the definition and normalization of \(J_p\), the admissible class of convex bodies \(K\), nor the hypotheses and concentration bound of Theorem~2 are included in the supplied record. Consequently, this is an explicitly incomplete formulation: it does not specify a numerical candidate bound, its dependence on \(n,p,\dim L\), or whether the concentration inequalities are strict. A fully quantified sharp-threshold statement requires the omitted theorem and its definitions.
\subsection{Short English statement}
The problem asks how much of a measure may lie on great subspheres while still guaranteeing that it is the \(L_p\) integral curvature measure of a convex body, for \(p\leq -1\). The original target is the optimality of a bound in Theorem~2, whose statement and supporting definitions are absent. The supplied open-status assessment is dated curation, not an independently verified conclusion.
\subsection{Sources}
\begin{itemize}
\item Supplied local registry record: permanent UnsolvedMath ID \(30004953\), source problem code \texttt{catalogue entry}, in the Oberwolfach Reports Open Problems collection. No individual proposer is identified; the proposed year is recorded as 2021.
\item Supplied report citation: \emph{Convex Geometry and its Applications (2022)}, Oberwolfach Reports, pp.~3187--3266, DOI \texttt{10.4171/owr/2021/59}. Supplied source URL and literature-source link labelled ``original OWR report'': \url{https://doi.org/10.4171/owr/2021/59}. This is the supplied authoritative reference; its relevant theorem text is not included in the record.
\item The supplied original passage identifies the optimality of the measure concentration bound in Theorem~2 as an open problem. The background attributes the extraction to workshop open-problem discussions documented in Oberwolfach Reports.
\item Literature assessment dated exactly \texttt{2026-08-22}: the supplied curation labels the problem open and reports that no primary source located by its cutoff establishes optimality. The listed source checked is the original OWR report at \url{https://doi.org/10.4171/owr/2021/59}. This assessment is supplied dated curation, not independently verified; the background explicitly requests independent verification of time-sensitive status and cited claims.
\item The registry supplies the statement-status label \texttt{corrected\_verified} and asserts that the clean statement was checked against the official report and preserves its hypotheses, notation, and target. These are supplied verification claims, not verification performed for this draft.
\end{itemize}
\end{document}
